\section{Unserious Meditations}

I used to think that books like Brave New World and 1984\footnote{\url{http://www.ritholtz.com/blog/2010/08/who-was-right-huxley-or-orwell/}} were intended to warn us about what may happen. Now I realize their purpose was to prepare us to welcome the brave new world with open arms and closed minds. The real irony are those vocal people think they are still free when they are fully immersed in it.

Suppose we eat government recommended food, read government approved news, receive a government education, etc. After all the answers are given, the problem of life still remains.

I watched part of a movie called Hysteria\footnote{\url{http://en.wikipedia.org/wiki/Hysteria_(2011_film)}} tonight. The heroine, who was a feminist, a social reformer, a free love advocate, was described as the most Christian woman in England. The doctor got rich by developing a personal massager to bring woman to orgasm. Once women are liberated, prostitutes find honorable work, and men learn how to sexually satisfy a woman, the problems of life still remain.

I really can't watch a movie anymore. It only works when the viewer is drawn in to the film and identifies with the characters. However, when you see through it, when you see actors reciting lines, when you understand how props are used and stories developed, when you read the subliminal message, then there is little more than flashing multi-colored lights.

The same applies to life. When there are only discrete and unrelated transient events, a man cannot be attracted by the glamour of life. The trick is to find the unity behind those appearances.

The Exoteric and Esoteric are not opposed but are intimately related. The exoteric perspective describes things from the outside and tries to prove things through logic and the rational mind. The esoteric, on the other hand, takes the inner view and understands things by intuition, which is closer to seeing than thinking. The philosopher is exoteric and the metaphysician is esoteric.

For example, regarding the soul, the exoteric will name its components and argue for its findings. The esoteric, instead, will observe the inner workings of the soul. The insoluble philosophical problem is the one and many. The metaphysician, through meditation, will see the one and the many directly. He will observe the ever changing contents of consciousness. He will learn to discern the different levels and will seek the unifying force behind all those appearances, that which does not change amidst all the change. That is itself not an appearance.

The exoterist will follow God's will in all things and trusts that God has chosen all for him. The esoterist will see how he chose his own life and all its circumstances. He does not mean anything different from the exoterist. Atman is Brahman, or put on the mind of Christ through whom all things were made.

Nietzsche criticized bourgeois Christians of his time, but he could not have criticized Christianity itself, which is not a ``thing''. A man doesn't choose a worldview because of its alleged consequences but because he believes it is true.

Nietzsche's critiques are not surprising to those who know, yet that seem novel to others. That true religion declines because of the degeneration of castes, is one such fact. That the priesthood becomes effeminate when inferior castes take dominion is another. The half-educated think they have discovered universal, rather than merely contingent, truths.

Children at school learn to ``show and tell''. Nietzsche tells, but he seldom shows.

That is because for Nietzsche, nothing is known until the act is performed. Specifically, thought follows action and, at best, it serves an instrumental purpose to make action effective. At worse, it is a form of false consciousness that hides the will to power.

Nietzsche predicted dire consequences because of Western Christianity. However, there have been ever worse consequences from his thought. The first is that logical discussion is no longer possible because of the hermeneutics of suspicion. All thought is presumed to be the ideological superstructure hiding the desire for money, sex and power; their respective prophets are Marx, Freud, and Nietzsche.

The second consequence is the idea that a man can become a superman simply by rejecting Christianity. Doing so does not make a man
\emph{ipso facto} a superman, no more than putting on a muscle shirt

makes a man a bodybuilder. Europe, for example, has been repaganized. Where, then, do we see a resurgence of Nietzschean supermen? But the neopagans would immediately object: they are not the real pagans, we are. So likewise the bourgeois Christians of the 19\textsuperscript{th} century are not the real Christians. So who are?

\textbf{Pope Leo XIII} tells us in \textit{Immortale Dei}\footnote{\url{http://www.vatican.va/holy_father/leo_xiii/encyclicals/documents/hf_l-xiii_enc_01111885_immortale-dei_en.html}}:

\begin{quotex}
There was once a time when States were governed by the philosophy of the Gospel. Then it was that the power and divine virtue of Christian wisdom had diffused itself throughout the laws, institutions, and morals of the people, permeating all ranks and relations of civil society. Then, too, the religion instituted by Jesus Christ, established firmly in befitting dignity, flourished everywhere, by the favour of princes and the legitimate protection of magistrates; and Church and State were happily united in concord and friendly interchange of good offices.

The State, constituted in this wise, bore fruits important beyond all expectation, whose remembrance is still, and always will be, in renown, witnessed to as they are by countless proofs which can never be blotted out or ever obscured by any craft of any enemies. Christian Europe has subdued barbarous nations, and changed them from a savage to a civilized condition, from superstition to true worship. It victoriously rolled back the tide of Mohammedan conquest; retained the headship of civilization; stood forth in the front rank as the leader and teacher of all, in every branch of national culture; bestowed on the world the gift of true and many-sided liberty; and most wisely founded very numerous institutions for the solace of human suffering. And if we inquire how it was able to bring about so altered a condition of things, the answer is beyond all question, in large measure, through religion, under whose auspices so many great undertakings were set on foot, through whose aid they were brought to completion.
\end{quotex}


The anti-Christians, then, must dispute that description or risk being irrelevant. They must refute the intellectual founders of the Middle Ages: \textbf{Plotinus}, \textbf{Augustine}, \textbf{Boethius}. They must show that \textbf{Thomas Aquinas} does not represent traditional metaphysics. They must deny that \textbf{Dante Alighieri} represents the peak of medieval Christian thought. Those who claim to be traditionalists must show how \textbf{Rene Guenon}, \textbf{Julius Evola}, and \textbf{Ananda Coomaraswamy} were wrong about the Middle Ages and tradition.

Dante regarded the Roman Empire as the supreme manifestation of temporal power and the Medieval church as the spiritual authority.

In ancient Rome, punishment was public for the edification of others; in the Middle Ages, punishment was usually hidden, since it was for the edification of the offender.

Only the isolated ascetics can restore tradition. Even if that is true, he would know how to recognize them.

Here is one such man. \textbf{St Mark the Ascetic} in the
\emph{Philokalia} mentions three great enemies:


\begin{enumerate}


\item Ignorance

\item Forgetfulness

\item Sloth
\end{enumerate}


\textbf{Ignorance}, as the enemy of the Intellect, is the first thing to overcome. Obviously a man must learn proper doctrine, not just exoteric, but also esoteric.

\textbf{Forgetfulness} results from being absorbed in the glamour of life. Consciousness identifies with the many rather than the one. Never forget who you are, why you were born, what you need to accomplish, what your real life is. Of course, most humans prefer to drink from Lethe, the river in Hell that causes them to forget.

\textbf{Sloth} is the enemy of the Will and the abandonment of the will to power. No matter how much you know, it needs to be brought into manifestation, that is, real world effects. Therefore, you will never really know what is possible, nor even who you really are, unless and until the right efforts are made.


\flrightit{Posted on 2013-02-25 by Cologero}

\begin{center}* * *\end{center}

\begin{footnotesize}
\begin{sffamily}
\texttt{Cassiodorus on 2013-02-26 at 00:25 said:}

``The Exoteric and Esoteric are not opposed but are intimately related.''

The teachings of Tradition would go on to say that doctrine, though indispensable, is \textbf{transcended} from \emph{within}. This is the key principle behind the Transcendent Unity of Religions.

At OrthodoxChristianity.net, I came across an interesting debate on this subject that I would like to share. In this exchange, a defender of Christian exclusivism insisted that Eastern Orthodoxy is, of course, fundamentally esoteric with Divinization as its central teaching. However, he argued against the Perennialist view that esoterism transcends theology and doctrine.

\emph{``Our (Orthodoxy) ``esoterism'' is inextricably bound up with the ``exoteric'' dogmas and ones leads to the other. Difference in dogma results in difference of spiritual experience. Our theology and worship are inherently mystical... ... those further on the path of deification will grow deeper in their understanding of Revelation''}

If the traditional Christian insists that his theology IS his esoterism, how are we to understand this with the Traditionalist framework?

\hfill

\texttt{Avery Morrow on 2013-02-26 at 02:03 said:}

Didn't Guénon say something like this with regards to Christianity?

\emph{The influence that operates through the medium of the Christian sacraments, having originally acted in the initiatic order,  subsequently, in other circumstances and for reasons dependent upon these,  lowered its action to the simply religious and exoteric domain;  this occurred in such a way that its effects have ever since been limited to certain possibilities of an exclusively individual order, whose scope does not extend beyond salvation... }

In other words, in agreement with your Orthodox friend, Christianity lacks a limited and initiatory core; once the sacraments became exoteric, they were no longer esoteric. I do not see anything in Guénon's claim which is much more controversial than this. But this drove Schuon batty; he replies in \emph{René Guénon: Some Observations} that \emph{we see absolutely no reason for this `lowering', in a case like that of Christianity, because the Spirit can in any event proportion its activity according to the capacity of the human receptacle; is God so poor that he would need to ration his graces after having granted them? Why should not one and the same rite be able to confer individual help to one, and supra-individual help to another? For he who can do the greater can do the lesser.}

I am not clear whether Schuon understands the need for limited initiation and the concept of caste which is so key to both Guénon and Evola; some events in his later life reflect this. The more Schuon I read, the less I understand him.

\hfill

\texttt{Mihai on 2013-02-26 at 04:39 said:}

A link to that specific topic would be useful.

I completely agree with his view regarding the esoteric-exoteric issue. I always say that, in Christianity, the degree of esoterism is directly proportional to the degree of inner understanding.

See Vladimir Lossky- The mystical theology of the Eastern Church (by the way, around here, the word ``mystical'' is used, rather than ``esoteric''. However, it is not the sentimental mysticism of the west, but it must be taken in the original meaning of the word, which is the same as ``esoteric''. )

Also, Cologero, I really identify with this line: ``I really can't watch a movie anymore. ``

\hfill

\texttt{Cassiodorus on 2013-02-26 at 12:04 said:}

Mihai,

Since I seem to lack the capacity to create a link, you can google the ``Theology of Professor James Cutsinger'' at OrthodoxChristianity.net. There is also some very interesting discussion on this subject to be found at Gnosisandnoesis.net- search James Cutsinger and/or Seyyed Hossain Nasr.

Avery,

I think Schuon had said, basically, that Christianity was an esoterism that was ``exteriorized'' to save the decaying Greco-Roman world- an ``eso-exorterism''?

I did see a movie recently that wasn't good, but it had something, I don't know... ..well, interesting? Valhalla Rising.

\hfill

\texttt{Saladin on 2013-02-26 at 19:49 said:}

``I really can't watch a movie anymore.'' Wow! I am not the only one!

\hfill

\texttt{Logres on 2013-02-27 at 10:40 said:}

I think for me that began with the Matrix, second part, when, after the orgy in Zion, I began to root for the Machines. Although Bergman's The Magician was interesting to see, lately.

\hfill

\texttt{Cassiodorus on 2013-02-27 at 16:12 said:}

I streamed an old classic last night off netflix. Definitely worth a re-watch if you haven't seen it- ``A Man For All Seasons''. The last scene where Sir Thomas Moore is on trial is excellent!!!

\end{sffamily}
\end{footnotesize}

